\chapter*{Conclusions}
\addcontentsline{toc}{chapter}{Conclusions}

This project has investigated the delicate balance between cost, usability, and security in the modern consumer Internet of Things (IoT) landscape, using the Shelly Plug S (Gen~1) as a concrete testing ground. By adopting both offensive and defensive perspectives, the research has traced an educational and engineering path that ranges from hardware reverse engineering to the deployment of cryptographic defenses on resource-constrained microcontrollers.

\medskip
\noindent The project journey began with the physical dissection of the device and the extraction of its binary firmware via the exposed UART interface. This initial phase highlighted a common paradigm in legacy consumer IoT: the reliance on cleartext protocols and unprotected storage to minimize development overhead and unit costs. The discovery of critical network credentials, cloud encryption keys, and device identifiers stored in plaintext within the flash memory, combined with the use of unencrypted MQTT and HTTP protocols, confirmed that the device lacked transport-layer and storage-level protection. 

\medskip
\noindent In the second phase, we translated these findings into a practical demonstration of risk by developing a custom firmware capable of silently exfiltrating telemetry data via UDP while maintaining behavioral camouflage via MQTT to prevent detection. The successful deployment of this firmware through network-level hijacking (ARP poisoning and DNS spoofing) and the direct abuse of the unauthenticated OTA endpoint demonstrated how easily an attacker within the local network can compromise device integrity. Crucially, the subsequent analysis of the exfiltrated power consumption profiles showed that the privacy risks associated with compromised smart plugs extend far beyond simple control interception, enabling the reconstruction of user habits, presence cycles, and appliance usage patterns through side-channel telemetry.

\medskip
\noindent The final stage of the project completed the engineering loop by implementing robust mitigation strategies directly on the target hardware. Despite the severe memory constraints of the ESP8266 SoC, we successfully migrated the system to MQTTS with mutual TLS (mTLS) authentication using ECC certificates, ensuring both transport confidentiality and strong bidirectional identity verification. Simultaneously, the implementation of an RSA-2048 signature verification scheme utilizing PKCS\#1 v1.5 padding and SHA-256 hashing successfully secured the OTA update interface, preventing any future execution of unauthorized code. Wireshark analysis validated the effectiveness of these interventions, showing that all previously cleartext telemetry data was rendered entirely opaque to unauthorized observers.

\medskip
\noindent Under an ethical perspective, this research highlights the importance of responsible disclosure practices and security awareness in the IoT ecosystem. Because the analyzed vulnerabilities affect a legacy, first-generation commercial product whose weaknesses are documented and largely mitigated in newer hardware iterations (such as the ESP32-based Gen~2 and Gen~3 lines), this work did not target unpatched zero-day vulnerabilities. Instead, the experimental activities were conducted within a fully isolated laboratory network on personally owned hardware, ensuring that no commercial cloud infrastructure or third-party users were impacted. In a broader context, this project reinforces the necessity for IoT manufacturers to coordinate with the cybersecurity research community, adopting transparent disclosure frameworks that protect consumers by ensuring patches are deployed before architectural vulnerabilities are publicly discussed.